% ============================================================================
\section{Experimental Protocol}
\label{sec:experiments}
% ============================================================================

The protocol first tests individual transfer, then separates harmful mean
updates from sampling failures at common checkpoints, and finally evaluates
prediction, capability, and computation. End-to-end measurements are not yet
available; the results below specify the evidence to collect.

\subsection{A four-domain distillation setting}

We study mathematics, code, instruction following, and science with a
Qwen3-1.7B student in non-thinking mode. Mathematics and instruction
following use domain-specific Qwen3-1.7B RL teachers. Code and science
are routed to the same frozen Qwen3-4B model through separate task
endpoints. This setting combines specialized feedback and a larger
generalist teacher in one student.

The default loss is Equation~\ref{eq:dense_topk_loss} with teacher
$k=64$, and each task has weight $w_i=1/4$. Each step processes
$G=16$ micro-batches of $b=4$ responses, with $2\leq m_i\leq8$.
Uniform uses four micro-batches per task. Runs last $500$ steps;
each task supplies a shuffled stream of $16{,}000$ unique prompts,
and each response is capped at $4{,}096$ tokens. There is one optimizer
update per fresh rollout batch. Noise and time estimates use EMA decay
$0.9$. Runs use AdamW and shared task prompt orders within each paired seed.
Uniform and GPAS (per step) use three paired training seeds; the other seven
joint configurations use one exploratory seed each.

One $96$-GB GPU trains the student, and one $48$-GB GPU serves
student rollouts and teacher scoring. The two generalist task endpoints
share their teacher weights in memory. Appendix~\ref{app:reference_tables}
records the configuration.
\missing{Complete the exact checkpoint and data revisions, optimizer
settings, software versions, and measured resource use.}

\subsection{Individual transfer and joint training}

Four single-task OPD references use the same initial student, assigned teacher,
and default loss with weight one. Each initially uses one seed, $500$ updates,
and $16$ micro-batches per update. Its $32{,}000$ responses require two passes
through the $16{,}000$-prompt pool, reshuffling only after exhaustion. Joint
runs do not repeat prompts. We report unique prompts, exposure, and compute,
and compare learning curves at matched domain exposure and total compute.
The reference's final score is an observed separately acquired capability,
neither an upper bound nor evidence of joint attainability.

The fixed-objective comparisons change only the assignment of the $16$
micro-batches:
\begin{itemize}
    \item \textbf{Uniform}: four micro-batches per task.
    \item \textbf{GPAS (per step)}: counts from smoothed preconditioned
    gradient noise.
    \item \textbf{GPAS (cost-aware)}: counts from the time--variance
    criterion in Equation~\ref{eq:batch_cost_objective}.
    \item \textbf{Unpreconditioned allocation}: the same variance rule
    with $D=I$ for count selection; AdamW still performs the update.
    \item \textbf{D$^3$ signal, fixed weights}: counts from D$^3$-MOPD's
    remaining-gap and descent-velocity signal, combined with our fixed
    task-mean weights.
\end{itemize}

Two recipe comparisons change the objectives. \textbf{D$^3$-MOPD scheduling}
averages responses across the step, giving effective task weights $m_i/G$.
\textbf{Open-MOPD ($K=1$)} uses its student top-$16$ dense loss, token-share
balancing, and gap-aware weighting. Appendix~\ref{app:baselines} specifies
these adaptations to our models, data, and budget.
\textbf{TA-OPD + Uniform} and \textbf{TA-OPD + GPAS} use the tail-aware loss
of \citet{huang2026tailaware} at $k=64$ to test loss sensitivity. These nine
joint configurations accompany the four single-task references. Only the
repeated Uniform--GPAS comparison supports claims about training-seed stability.

\subsection{Interference and precision at common checkpoints}

At Uniform steps $50$, $250$, and $450$, each diagnostic branch restores the
same complete student, AdamW, scheduler, data, and random state and makes one
update. We compare each task alone, the Uniform joint update, GPAS, and a
more precise equal-weight joint estimate using $32$ micro-batches per task.
Single-task branches use $16$ micro-batches of their task with weight one;
Uniform and GPAS retain $G=16$. The precise branch is an expensive diagnostic,
and its cost is reported separately.

At each checkpoint, calibration, update-trial, and evaluation samples are
drawn independently from the same declared diagnostic distribution for each
task, using the fixed checkpoint rollout policy. These distributions are
held out from training and benchmarks. For these tests, $L_i^t$ and $\mu_i$
refer to the diagnostic distribution; applicability to the training
distribution is evaluated empirically. Draws and responses are not reused
across roles; independent draws may select the same prompt. Calibration uses
$32$ independent micro-batches per task to estimate means, trace noise,
transfer margins $a_i$ (Equation~\ref{eq:transfer_margins}), and directional
variances (Equation~\ref{eq:directional_variance}). We split calibration into
two independent halves and use cross-fitted mean products for $K$ and $a_i$,
avoiding the upward bias from squaring the same noisy mean estimate.
Calibration estimates alone select diagnostic GPAS counts. We run $20$
independent update trials per branch and evaluate every branch on the same independent, fixed reference
prefixes for all four domains. Diagnostic samples are independent with replacement; joint training uses
finite streams without replacement. Their covariance assumptions are distinct.

For each trial we use independent evaluation gradients to estimate the
first-order projection under frozen pre-step $D$, and record the fixed-prefix
loss change after an actual AdamW update. A separate equal-allocation sweep uses
$m_i\in\{2,4,8,16,32\}$ for every task, so $G$ ranges from $8$ to $128$.
This local sweep is exempt from the main count cap and does not test the
end-to-end effect of changing update frequency under a fixed compute budget.

We classify a harmful mean interaction only when an independently evaluated
margin estimate is negative beyond its sampling uncertainty. A precise AdamW
update can increase loss because of momentum or finite-step effects even with
a positive frozen-$D$ margin; those cases test the approximation. Positive
margins with frequent sampled sign reversals identify a candidate sampling
failure. Margins whose uncertainty includes zero remain inconclusive.
Total trace noise is assessed as a proxy for directional uncertainty; a directional allocation
is considered only as an offline diagnostic if that proxy fails.

\subsection{Prediction before intervention}

We specify the diagnostic rule, regression threshold, confidence procedure,
and comparison metrics using step $50$ of the first Uniform seed, before
examining later diagnostic outcomes. The rule uses estimated margin signs and
$v_i/a_i^2$ for positive margins to predict domain regressions, the reduction
in regressions from GPAS or increased precision, and cases with little expected
benefit. Steps $250$ and $450$ and the two other training seeds evaluate these
predictions without retuning the rule. We compare predictions from directional
uncertainty with those from trace noise alone, and report predictive errors
as well as successful interventions. Plug-in estimates of the probability
bound are diagnostic quantities, not finite-sample certificates.
\missing{Record the rule and analysis specification before evaluating the
held-out checkpoints and seeds.}

\subsection{Capability, state distribution, and computation}

\paragraph{Separate measurements.}
We report MATH-500 greedy pass@1, a fixed LiveCodeBench pass@1 slice,
IFBench strict accuracy, and GPQA-Diamond average@4. The initial student and
assigned teachers use the same benchmark protocols. Raw per-domain scores
remain primary. Reference-normalized gains are reported only when the
single-task gain exceeds a prespecified positive denominator floor; otherwise
we report raw gains. Every $50$ steps, all methods are evaluated with the
common teacher top-$64$ loss and equal task weights on $128$ fixed held-out
prompts per task. We keep two distinct loss measurements: a fixed bank of
initial-student reference prefixes, and fresh prefixes generated by the
current student. Common-checkpoint diagnostics additionally use an independent
bank generated at that checkpoint and held fixed across its branches.
We retain each checkpoint's fresh reference bank and teacher scores, then
evaluate the next checkpoint on both the previous and new banks. These paired
measurements give the empirical counterpart of Equation~\ref{eq:policy_drift}
on the held-out prompt distribution over $50$-step intervals, with sampling
uncertainty reported separately. The initial bank provides a common fixed-loss target;
its difference from fresh-policy loss alone is not that decomposition.
Neither loss is treated as a capability score.

The main capability table uses step $500$ for every run, without selecting a
different checkpoint for each domain. For the repeated core comparison,
we report all three seed outcomes and their mean and spread. Paired
prompt-level bootstrap intervals describe evaluation uncertainty separately;
they do not replace training-seed variation. Single-task references initially
have one seed, which limits claims about their stability.
\missing{Insert benchmark versions, code evaluation date range, complete
decoding settings, the reference-gain denominator floor, and capability
tolerances $\epsilon_i$ before inspecting results.}

\paragraph{Measured efficiency.}
We plot fixed-reference loss, fresh-policy loss, and capability separately
against task exposure, optimizer steps, and GPU hours. GPU hours include all
time on the two reserved devices, including waiting, synchronization,
statistics, and checkpointing. Diagnostic branches and evaluation have
separate compute totals. Within the fixed-weight default-loss group, we
report time to Uniform's mean final fixed-reference loss, interpolating
between adjacent evaluations and marking unreached targets. This is a
loss-based efficiency measure; per-domain capability and the worst-domain
change accompany it.

\paragraph{Allocation traces.}
The training loop logs task loss, its recent rate of change, response length,
trace noise with and without preconditioning, assigned counts, task exposure,
and time per micro-batch. Diagnostic checkpoints add progress margins and
directional noise. We test whether these measurements explain observed
regressions and intervention effects. The cost-aware criterion measures time
times update variance; its relationship to learning progress is evaluated
empirically. All statistics collection costs enter the time and memory report.

% ============================================================================
\section{Results}
\label{sec:results}
% ============================================================================

\subsection{Separately acquired and jointly acquired capabilities}

\begin{figure}[t]
    \centering
    \missingfigure[0.17\textheight]{Plot single-task and joint capability
    gains per domain against domain exposure and total GPU hours. Mark prompt
    repetition in single-task references and show raw scores. Measurements
    are not yet available.}
    \caption{Individual transfer and joint training under explicit exposure
    and compute comparisons. A difference does not identify its cause.}
    \label{fig:individual_joint}
\end{figure}

Figure~\ref{fig:individual_joint} tests whether the chosen student acquires
each capability separately and whether joint training reaches comparable
gains. A task that does not improve alone does not establish an integration
failure. A joint-versus-single-task gap requires the common-checkpoint tests
before it can be attributed to interference or sampling.
\missing{Report raw gains, uncertainty, matched-exposure comparisons, and
any tasks for which individual transfer is weak or inconclusive.}

\subsection{Mean interactions, sampling failures, and prediction}

\begin{figure}[t]
    \centering
    \missingfigure[0.19\textheight]{For each checkpoint and domain, show
    cross-fitted progress margins with uncertainty, directional and trace noise,
    independent first-order sign measurements, and actual fixed-prefix loss changes for
    single-task, Uniform, GPAS, and precise-gradient branches. Add the local
    batch-size sweep. All panels await measurements.}
    \caption{Common-checkpoint interventions distinguish estimated mean
    interactions from sensitivity to sampling and test the frozen-$D$
    approximation against actual AdamW updates.}
    \label{fig:checkpoint_diagnostics}
\end{figure}

\missing{Report which estimated mean margins are positive, negative, or
inconclusive; where sampled projections disagree with them; and whether
increasing precision reduces the observed failures. Include uncertainty in the precise
estimate and cases where finite-step AdamW behavior differs from the local
projection. Then report predictive performance on held-out checkpoints and
seeds, comparing directional uncertainty with trace noise. Do not infer
capacity limits or long-run impossibility from local regressions.}

\subsection{Capability and computation after intervention}

\begin{table}[t]
\centering
\small
\setlength{\tabcolsep}{4pt}
\caption{Four-domain capability at step $500$. Blank cells await measurements.
Uniform and GPAS use three seeds; all other joint configurations and each
single-task reference initially use one. The single-task row contains four
separate students, each scored in its training domain, not one joint model.}
\label{tab:main_results}
\begin{tabular}{lccccc}
\toprule
Method & Math & Code & IF & Science & Mean \\
\midrule
Initial student & & & & & \\
Assigned teachers & & & & & \\
Single-task references & & & & & \\
\midrule
Uniform & & & & & \\
GPAS (per step) & & & & & \\
\midrule
GPAS (cost-aware) & & & & & \\
Unpreconditioned allocation & & & & & \\
D$^3$ signal, fixed weights & & & & & \\
D$^3$-MOPD scheduling & & & & & \\
Open-MOPD ($K=1$) & & & & & \\
TA-OPD + Uniform & & & & & \\
TA-OPD + GPAS & & & & & \\
\bottomrule
\end{tabular}
\end{table}

\begin{figure}[t]
    \centering
    \missingfigure[0.17\textheight]{Plot fixed-reference loss, fresh-policy
    loss, and per-domain capability in separate panels against steps and
    measured GPU hours. Show three-seed variation for Uniform and GPAS;
    label secondary comparisons as single-seed exploratory runs.
    Measurements are not yet available.}
    \caption{Training progress, changes in visited states, and capability
    evaluated separately under the same compute accounting.}
    \label{fig:training_efficiency}
\end{figure}

Table~\ref{tab:main_results} and Figure~\ref{fig:training_efficiency} test
whether the diagnosed local effect leads to useful end-to-end changes.
\missing{Report every per-domain change, the worst-domain change, seed
variation, measured GPU hours, and time to the fixed-reference loss target.
State whether capability and the two loss measurements agree. Report negative
or null results alongside improvements.}

\subsection{Allocation, cost, and dense-loss sensitivity}

\begin{figure}[t]
    \centering
    \missingfigure[0.16\textheight]{Align task counts, trace noise, loss,
    and response length over training; mark the diagnostic checkpoints and
    their progress margins. Keep preconditioned and unpreconditioned noise
    on separate scales. Measurements are not yet available.}
    \caption{Allocation traces interpreted using the measured progress
    margins and intervention outcomes.}
    \label{fig:allocation_dynamics}
\end{figure}

The unpreconditioned and fixed-weight D$^3$ comparisons test the role of
optimizer scaling and the choice of allocation signal. The cost-aware run
tests whether minimizing time times variance improves measured efficiency.
The TA-OPD pair tests sensitivity to the dense loss.
\missing{Relate allocation changes to measured directional uncertainty and
outcomes. Report statistics overhead, fitted and observed task costs, and
capability changes under the second loss. Keep conclusions from single-seed
comparisons exploratory.}
